%
% Academic CV --- Taejun Ha
% Based on a template by Dubasi Pavan Kumar (MIT License)
% Restructured for PhD applications
%

\documentclass[a4paper,11pt]{article}

% Package imports
\usepackage{latexsym}
\usepackage{xcolor}
\usepackage{float}
\usepackage{ragged2e}
\usepackage[empty]{fullpage}
\usepackage{wrapfig}
\usepackage{tabularx}
\usepackage{titlesec}
\usepackage{geometry}
\usepackage{marvosym}
\usepackage{verbatim}
\usepackage{enumitem}
\usepackage{fancyhdr}
\usepackage{multicol}
\usepackage{graphicx}
\usepackage{cfr-lm}
\usepackage[T1]{fontenc}
\usepackage{fontawesome5}
\usepackage[hidelinks]{hyperref}
\usepackage[most]{tcolorbox}

% Color definitions
\definecolor{darkblue}{RGB}{0,0,139}

% Page layout
\setlength{\multicolsep}{0pt}
\geometry{
    left=1.5cm,
    top=1.5cm,
    right=1.4cm,
    bottom=1.5cm
}

% Running header: name appears on page 2 onward
\pagestyle{fancy}
\fancyhf{}
\fancyhead[L]{\textit{\small Taejun Ha}}
\fancyhead[R]{\textit{\small Curriculum Vitae}}
\renewcommand{\headrulewidth}{0pt}
\renewcommand{\footrulewidth}{0pt}
\setlength{\headheight}{14pt}
\setlength{\footskip}{5pt}

% Hyperlink setup
\hypersetup{
    colorlinks=true,
    linkcolor=darkblue,
    filecolor=darkblue,
    urlcolor=darkblue
}

% Custom box settings
\tcbset{
    frame code={},
    center title,
    left=0pt,
    right=0pt,
    top=0pt,
    bottom=0pt,
    colback=gray!20,
    colframe=white,
    width=\dimexpr\textwidth\relax,
    enlarge left by=-2mm,
    boxsep=4pt,
    arc=0pt,
    outer arc=0pt
}

% URL style
\urlstyle{same}

% Text alignment
\raggedright
\setlength{\tabcolsep}{0in}

% Section formatting
\titleformat{\section}{
    \vspace{2pt}\scshape\raggedright\large
}{}{0em}{}[\color{black}\titlerule \vspace{-2pt}]

% Custom commands
\newcommand{\resumeItem}[2]{
    \item{
        \textbf{#1}{\hspace{0.5mm}#2 \vspace{0.3mm}}
    }
}

\newcommand{\resumePOR}[3]{
    \vspace{0.5mm}\item
    \begin{tabular*}{0.97\textwidth}[t]{l@{\extracolsep{\fill}}r}
        \textbf{#1}\hspace{0.3mm}#2 & \textit{\small{#3}}
    \end{tabular*}
    \vspace{-2mm}
}

\newcommand{\resumeSubheading}[4]{
    \vspace{0.5mm}\item
    \begin{tabular*}{0.98\textwidth}[t]{l@{\extracolsep{\fill}}r}
        \textbf{#1} & \textit{\footnotesize{#4}} \\
        \textit{\footnotesize{#3}} & \footnotesize{#2} \\
    \end{tabular*}
    \vspace{-1.2mm}
}

\newcommand{\resumeProject}[4]{
    \vspace{0.5mm}\item
    \begin{tabular*}{0.98\textwidth}[t]{l@{\extracolsep{\fill}}r}
        \textbf{#1} & \textit{\footnotesize{#3}} \\
        \footnotesize{\textit{#2}} & \footnotesize{#4}
    \end{tabular*}
    \vspace{-2.4mm}
}

\newcommand{\resumeSubItem}[2]{
    \resumeItem{#1}{#2}\vspace{-4pt}
}

\renewcommand{\labelitemi}{$\vcenter{\hbox{\tiny$\bullet$}}$}
\renewcommand{\labelitemii}{$\vcenter{\hbox{\tiny$\circ$}}$}

\newcommand{\resumeSubHeadingListStart}{
    \begin{itemize}[leftmargin=*,labelsep=1mm]
}

\newcommand{\resumeHeadingSkillStart}{
    \begin{itemize}[leftmargin=*,itemsep=1.7mm,rightmargin=2ex]
}

\newcommand{\resumeItemListStart}{
    \begin{itemize}[leftmargin=*,labelsep=1mm,itemsep=1.4mm]
}

\newcommand{\resumeSubHeadingListEnd}{
    \end{itemize}\vspace{2mm}
}

\newcommand{\resumeHeadingSkillEnd}{
    \end{itemize}\vspace{-2mm}
}

\newcommand{\resumeItemListEnd}{
    \end{itemize}\vspace{-0.5mm}
}

\newcolumntype{L}{>{\raggedright\arraybackslash}X}
\newcolumntype{R}{>{\raggedleft\arraybackslash}X}
\newcolumntype{C}{>{\centering\arraybackslash}X}

% Commands for icon sizing and positioning
\newcommand{\socialicon}[1]{
    \raisebox{-0.05em}{\resizebox{!}{1em}{#1}}
}

% Font options
\newcommand{\headerfonti}{\fontfamily{phv}\selectfont}
\newcommand{\headerfontii}{\fontfamily{ptm}\selectfont}
\newcommand{\headerfontiii}{\fontfamily{ppl}\selectfont}
\newcommand{\headerfontiv}{\fontfamily{pbk}\selectfont}
\newcommand{\headerfontv}{\fontfamily{pag}\selectfont}
\newcommand{\headerfontvi}{\fontfamily{cmss}\selectfont}
\newcommand{\headerfontvii}{\fontfamily{qhv}\selectfont}
\newcommand{\headerfontviii}{\fontfamily{qpl}\selectfont}
\newcommand{\headerfontix}{\fontfamily{qtm}\selectfont}
\newcommand{\headerfontx}{\fontfamily{bch}\selectfont}

\begin{document}

\headerfontiii
\thispagestyle{empty}

% ============================================================
% HEADER
% ============================================================

\begin{center}
    {\Huge\textbf{TAEJUN HA}}
\end{center}
\vspace{-2mm}

\begin{center}
    \small{\href{mailto:tjha@yonsei.ac.kr}{tjha@yonsei.ac.kr}}
\end{center}
\vspace{-5mm}

\begin{center}
    \small{
        \socialicon{\faGraduationCap}
        \href{https://scholar.google.com/citations?user=SCbCRH0AAAAJ}{Google Scholar}
        $\vert$
        \socialicon{\faGithub}
        \href{https://github.com/TaeJun1999}{GitHub}
        $\vert$
        \socialicon{\faOrcid}
        \href{https://orcid.org/0009-0004-8856-996X}{ORCID}
    }
\end{center}
\vspace{-3mm}

% ============================================================
% RESEARCH INTERESTS
% ============================================================

\section{\textbf{Research Interests}}
\vspace{1mm}

\small{
Machine learning for physical-layer wireless communications: model-based deep learning for
channel estimation, attention transformers and generative models (e.g.\ diffusion) as learned
signal priors, joint channel estimation and decoding, and low-complexity receiver design under
limited pilots and imperfect training labels. Broader interests include MIMO transceiver
optimization and physical-layer design for AI-native 6G and virtualized RAN.
}

\vspace{-2mm}

% ============================================================
% EDUCATION
% ============================================================

\section{\textbf{Education}}
\vspace{-0.4mm}

\resumeSubHeadingListStart

\resumeSubheading
{Yonsei University}
{Seoul, South Korea}
{M.S. in Electrical and Electronic Engineering (expected)}
{Mar 2025 -- Feb 2027}

\resumeItemListStart
\item Advisor: Prof. Jeonghun Park
\item Thesis: Learning Linear Filters for OFDM Channel Estimation Using an Attention Transformer
\resumeItemListEnd

\resumeSubheading
{Yonsei University}
{Seoul, South Korea}
{B.S. in Electrical and Electronic Engineering}
{Mar 2019 -- Feb 2025}



\resumeSubHeadingListEnd
\vspace{-4mm}

% ============================================================
% PUBLICATIONS AND PATENTS
% ============================================================

\section{
    \textbf{Publications and Patents}
    \hfill
    \textcolor{darkblue}{
        \scriptsize J=Journal, C=Conference, P=Patent
    }
}

\vspace{0.2mm}

\small{
\begin{enumerate}[
    leftmargin=*,
    labelsep=0.5em,
    align=left,
    widest={[\textbf{P.1}]},
    itemindent=0em,
    label={},
    itemsep=1.4mm
]

\item[\textbf{[J.1]}]
\textbf{T. Ha}, H. Kim, and J. Park,
``Attention-Aided MMSE with Ridge Denoising: How to Train under Noisy Channel Samples,''
\textit{IEEE Transactions on Vehicular Technology},
accepted for publication, 2026.
[\href{https://arxiv.org/abs/2608.15945}{Arxiv}]

\item[\textbf{[J.2]}]
\textbf{T. Ha}, C. Jung, H. Kim, J. Park, and J. Park,
``Learning MMSE Filters for OFDM Channel Estimation: Attention Transformer Gains at Linear
Inference,'' submitted to \textit{IEEE Transactions on Wireless Communications}.
[\href{https://arxiv.org/abs/2506.00452}{Arxiv}]
(Under Review, Minor Revision)

\item[\textbf{[J.3]}]
\textbf{T. Ha} and J. Park,
``Learning Linear Filters for Underwater OFDM Channel Estimation with Attention-Aided MMSE,''
\textit{Journal of the Korean Institute of Communications and Information Sciences (Letters)},
2026. (Published)

\item[\textbf{[C.1]}]
\textbf{T. Ha} and J. Park,
``Attention-Aided MMSE for OFDM Channel Estimation,''
in \textit{Proc. 59th Asilomar Conf. Signals, Systems, and Computers},
Pacific Grove, CA, USA, 2025.

\vspace{1mm}

\item[\textbf{[P.1]}]
J. Park and \textbf{T. Ha},
``Channel Estimation System and Method Using A-MMSE in OFDM Systems,''
Korean Patent Application No.~10-2025-0083228,
filed Jun.~24, 2025.
Applicant: Defense Acquisition Program Administration, Republic of Korea. (pending)

\item[\textbf{[P.2]}]
J. Park and \textbf{T. Ha},
``Method and Apparatus for Channel Estimation Using an Adaptive Channel-Estimation Filter,''
Korean Patent Application No.~10-2026-0127455,
filed Jul.~10, 2026.
Applicant: Yonsei University Industry-Academic Cooperation Foundation. (pending)

\end{enumerate}
}

% ============================================================
% RESEARCH EXPERIENCE
% ============================================================

\section{\textbf{Research Experience}}
\vspace{-0.4mm}

\resumeSubHeadingListStart

\resumeSubheading
{The University of Texas at Austin}
{Austin, TX, USA}
{Visiting Scholar, Dept.\ of Electrical and Computer Engineering}
{Jul 2026 -- Dec 2026}

\resumeItemListStart

\item Host: Prof. Hyeji Kim

\item Research topic: Diffusion-based joint channel estimation and decoding for
pilot-limited MIMO.

\item Method: Learned diffusion prior for the channel, BCJR decoder as an exact
symbol denoiser, coupled by extrinsic message passing.

\item Developed a ridge-regularized training objective and covariance-shrinkage surrogate
targets that allow attention-based channel estimators to be trained without clean CSI
labels, closing most of the gap to clean-label training at high SNR [J.1].

\resumeItemListEnd

\resumeSubheading
{Yonsei University}
{Seoul, South Korea}
{Graduate Research Assistant, Advisor: Prof. Jeonghun Park}
{Mar 2025 -- Present}

\resumeItemListStart

\item
\textbf{6GARROW: 6G AI-Native Integrated RAN-Core Networks}
(EU Horizon Europe SNS JU, Grant No.~101192194; with IITP, Korea) ---
Designed an attention-transformer channel estimator for 6G OFDM receivers that attains
near-MMSE accuracy at substantially lower computational cost, without requiring full
second-order channel statistics.

\item
\textbf{Development of Intelligent Secure Underwater Communication Technology}
(funded by KRIT, Korea) ---
Extended the attention-transformer estimator to underwater acoustic OFDM channels,
improving estimation accuracy and link reliability under long delay spread and severe
Doppler, supporting secure and dependable underwater links.

\item
Formulated the attention-aided MMSE (A-MMSE) framework, which learns linear
channel-estimation filters via an Attention Transformer and performs inference
through a single linear operation, removing the need for exact channel covariance;
developed a two-stage Attention encoder and a rank-adaptive extension for a flexible
performance-complexity trade-off. Presented at \textbf{Asilomar 2025} [C.1] and
extended into a journal version \textbf{under minor revision at IEEE Trans.
Wireless Communications} [J.2].

\resumeItemListEnd

\resumeSubHeadingListEnd
\vspace{-4mm}

% ============================================================
% GRADUATE COURSEWORK
% ============================================================

\section{\textbf{Graduate Coursework}}
\vspace{1mm}

\small{
Advanced Multi-Antenna System Design and Analysis,
Numerical Optimization,
Information Theory,
Broadband Wireless Communication Systems,
Random Processes,
Computer Vision,
Network System Analysis,
Advanced Programming.
}

\vspace{-2mm}

% ============================================================
% TECHNICAL SKILLS
% ============================================================

\section{\textbf{Technical Skills}}
\vspace{-0.4mm}

\resumeHeadingSkillStart

\resumeSubItem{Programming:}
{Python, MATLAB}

\resumeSubItem{Frameworks \& Tools:}
{PyTorch, NumPy/SciPy, Git, \LaTeX{}}

\resumeSubItem{Domain-Specific:}
{OFDM/MIMO link-level simulation; standardized channel models (COST 2100, 3GPP TDL/CDL);
LMMSE and model-based channel estimation; attention transformers and diffusion models
for physical-layer design; joint detection and decoding (BCJR, message passing)}

\resumeHeadingSkillEnd
\vspace{-2mm}

% ============================================================
% HONORS AND AWARDS
% ============================================================

\section{\textbf{Honors and Awards}}
\vspace{-0.4mm}

\resumeSubHeadingListStart

\resumePOR
{BK21 FOUR Graduate Fellowship}
{, Ministry of Education / NRF, Korea}
{2025}

\resumePOR
{DB Scholarship for Outstanding Incoming Students}
{, Yonsei University}
{2025}

\resumePOR
{Jinri Academic Excellence Scholarship}
{, Yonsei University (3 semesters)}
{2022 -- 2024}

\resumeSubHeadingListEnd
\vspace{-4mm}

% ============================================================
% TEACHING EXPERIENCE
% ============================================================

%\section{\textbf{Teaching Experience}}
%\vspace{-0.4mm}
%
%\resumeSubHeadingListStart
%
%\resumePOR
%{Teaching Assistant}
%{, [Course Name], Yonsei University}
%{[Semester Year]}
%
%\resumeSubHeadingListEnd
%\vspace{-2mm}

\end{document}